\documentclass[a4paper, 10pt]{article}

% --- PACOTES ---
\usepackage[utf8]{inputenc}
\usepackage[T1]{fontenc}
\usepackage[english,brazil]{babel}
\usepackage{geometry}
\usepackage{parskip}
\usepackage{hyperref}
\usepackage{xurl}
\usepackage{titlesec}
\usepackage{textcomp}
%\usepackage{enumitem}

% --- CONFIGS ---
\geometry{top=1.5cm, bottom=1.5cm, left=1.5cm, right=1.5cm}
\pagestyle{empty}

\hypersetup{
    pdftitle={Currículo - Wesley Roberto dos Santos},
    pdfauthor={Wesley Roberto dos Santos},
    colorlinks=true,
    linkcolor=black,
    urlcolor=black,
    citecolor=black
}

\setcounter{secnumdepth}{0}

% Formato da linha abaixo das seções
\titleformat{\section}
{\Large\bfseries}
{}
{0em}
{}
[\titlerule\vspace{0.5ex}]

% Listas compactas
%\setlist[itemize]{noitemsep, topsep=0pt}

\begin{document}

% --- Cabeçalho ---
\begin{center}
    {\Huge \textbf{Wesley Roberto dos Santos}} \\
    \vspace{0.3em}
    \small{Desenvolvedor Full Stack Júnior com foco em JavaScript (Node.js e React), com experiência prática na construção de aplicações web e integração com APIs.}

    \vspace{0.6em}
    Curitiba, PR, Brasil \quad\textbullet\quad 
    \href{mailto:wesleysantos32892653@outlook.com}{wesleysantos32892653@outlook.com} \quad\textbullet\quad
    LinkedIn: \href{https://www.linkedin.com/in/wesley-santos-04aa50122}{linkedin.com/in/wesley-santos-04aa50122} \quad\textbullet\quad
    GitHub: \href{https://github.com/TitanCodeXD}{github.com/TitanCodeXD} \quad\textbullet\quad
    Portfólio: \href{https://portfolio-wesley-santos.netlify.app}{portfolio-wesley-santos.netlify.app}
\end{center}

% --- Projetos ---
\section{Projetos}
\textbf{Aplicação Full-Stack no Estilo Instagram} \hfill 2024 \\
 Repositório no GitHub - \href{https://github.com/TitanCodeXD/Parallel}{https://github.com/TitanCodeXD/Parallel}
\begin{itemize}
    \item Desenvolvi uma aplicação web full-stack semelhante ao Instagram usando React.js, Node.js, MongoDB e Redux.
    \item Implementei funcionalidades como postagem de fotos, comentários, curtidas, upload de arquivos (AWS S3) e autenticação.
\end{itemize}

\medskip
\textbf{Aplicação Full-Stack de Chat} \hfill 2024 \\
Repositório no GitHub - \href{https://github.com/40nados/ProjectHub}{https://github.com/40nados/ProjectHub}
\begin{itemize}
    \item Criei uma plataforma de chat em tempo real com Next.js, Node.js, Socket.IO e MongoDB.
    \item Adicionei funcionalidades de upload de arquivos (AWS S3), autenticação e mensagens privadas/em grupo.
\end{itemize}

\medskip
\textbf{ChatBot Gemini} \hfill 2024 \\
 Repositório no GitHub - \href{https://github.com/TitanCodeXD/Chat-Bot}{https://github.com/TitanCodeXD/Chat-Bot}
\begin{itemize}
    \item Desenvolvi um chatbot com Google Gemini (IA generativa)
    \item Implementei uma interface interativa com Streamlit - Python, focada em usabilidade e experiência do usuário
    \item Fiz processamento de arquivos PDF, imagem e áudio com extração de informações de dados (BytesIO).
\end{itemize}

% --- Habilidades ---
\section{Habilidades e Certificações}
\begin{itemize}
    \item \textbf{Linguagens:} JavaScript, TypeScript, Python, C, R;
    \item \textbf{Frontend:} React.js, Next.js, HTML, CSS, Tailwind CSS;
    \item \textbf{Backend:} Node.js, NestJS, Express.js, MongoDB, SQL (SQLite, PostgreSQL), API RESTFul, Prisma, JWT;
    \item \textbf{Testes e Documentação:} Jest, Swagger/OpenAPI;
    \item \textbf{Ferramentas:} Git, Bash script, Docker, AWS S3, Redis, Firebase;
    \item \textbf{Idiomas:} Português (Nativo), Inglês (Intermediário), Russo (Básico);
    \item \textbf{Certificações:} MongoDB Associate Developer (Node.js) (2026), SCRUM Fundamentals Certified (2021);
    \item \textbf{Cursos:} FreeCodeCamp: JavaScript Algorithms and Data Structures (2021), Responsive Web Design (2021); Udemy: Node.js, Express, MongoDB, React e Hooks (2024).
\end{itemize}


% --- Formação ---
\section{Formação Acadêmica}
\textbf{Bacharelado em Engenharia da Computação} \hfill Ago 2019 -- Dez 2024 \\
Centro Universitário Tecnológico de Curitiba (UNIFATECPR), Curitiba, PR

\medskip
\textbf{Curso Técnico em Programação CNC e Operação de Torno} \hfill Jan 2018 -- Dez 2019 \\
SENAI, Curitiba, PR

% --- Experiência ---
\section{Experiência Profissional}
\textbf{Aprendiz de Mecânica e Auxiliar de Programação CNC} \hfill Mar 2019 -- Nov 2019 \\
Trombini Embalagens SA, Curitiba, PR
\begin{itemize}
    \item Colaborei na manutenção eficiente de máquinas, reduzindo o tempo total de inatividade em até 15\%.
    \item Auxiliei nos processos de programação CNC, aumentando efetividade nos códigos em até 10\%.
\end{itemize}

\end{document}
